% Projet Menia -- le programme de recherche explique simplement
% Compiler : pdflatex main.tex (deux fois, pour le sommaire)
\documentclass[11pt,a4paper]{article}

\usepackage[T1]{fontenc}
\usepackage[utf8]{inputenc}
\usepackage[french]{babel}
\usepackage[sc]{mathpazo}
\usepackage[scaled=0.92]{helvet}
\usepackage{microtype}
\usepackage{icomma}
\usepackage[a4paper,left=2.3cm,right=2.3cm,top=2.2cm,bottom=2.5cm]{geometry}
\usepackage{amsmath,amssymb}
\usepackage{xcolor}
\usepackage{tikz}
\usetikzlibrary{arrows.meta,positioning,calc,shapes.geometric,decorations.pathreplacing,babel,fit,backgrounds}
\usepackage{pgfplots}
\pgfplotsset{compat=1.18}
\pgfplotsset{/pgf/number format/use comma}
\usepackage[most]{tcolorbox}
\usepackage{booktabs,array,tabularx}
\usepackage{enumitem}
\usepackage{multicol}
\usepackage{titlesec}
\usepackage{fancyhdr}
\usepackage{caption}
\usepackage[hidelinks]{hyperref}

% ---------- Couleurs ----------
\definecolor{mbleu}{HTML}{1F5A96}
\definecolor{mvert}{HTML}{2E7D4F}
\definecolor{morange}{HTML}{C8641E}
\definecolor{mrouge}{HTML}{B03A2E}
\definecolor{mgris}{HTML}{5F6B73}
\definecolor{mjaune}{HTML}{D9A400}
\definecolor{mviolet}{HTML}{6A4C93}

% ---------- Mise en page ----------
\linespread{1.05}
\setlength{\parindent}{0pt}
\setlength{\parskip}{0.45em}
\setlist{itemsep=0.15em,topsep=0.3em,leftmargin=1.4em}
\captionsetup{font=small,labelfont={bf,sf},labelsep=period}

\titleformat{\section}{\sffamily\Large\bfseries\color{mbleu}}{\thesection\enspace·}{0.5em}{}
\titleformat{\subsection}{\sffamily\large\bfseries\color{mbleu!85!black}}{\thesubsection}{0.6em}{}
\titleformat{\subsubsection}{\sffamily\normalsize\bfseries\color{mgris}}{}{0em}{}
\titlespacing*{\section}{0pt}{1.4em}{0.6em}
\titlespacing*{\subsection}{0pt}{1.0em}{0.3em}

\pagestyle{fancy}
\fancyhf{}
\fancyhead[L]{\sffamily\small\color{mgris}Projet Menia}
\fancyhead[R]{\sffamily\small\color{mgris}expliqué simplement}
\fancyfoot[C]{\sffamily\small\thepage}
\renewcommand{\headrulewidth}{0.3pt}

% ---------- Commandes ----------
\newcommand{\q}[1]{\og #1\fg{}}
\newcommand{\pc}{\,\%}
\newcommand{\vect}[2]{(#1\,;\,#2)}

% ---------- Encadres ----------
\tcbset{enhanced,breakable,boxrule=0.6pt,arc=2mm,left=3mm,right=3mm,top=1.5mm,bottom=1.5mm,
  fonttitle=\sffamily\bfseries,before skip=0.7em,after skip=0.7em}
\newtcolorbox{retenir}{colback=mvert!6,colframe=mvert,title={À retenir}}
\newtcolorbox{exemple}[1][]{colback=mbleu!5,colframe=mbleu,title={Exemple à refaire à la main#1}}
\newtcolorbox{avous}{colback=morange!6,colframe=morange,title={À vous}}
\newtcolorbox{attention}{colback=mrouge!5,colframe=mrouge,title={Attention}}
\newtcolorbox{encours}{colback=mjaune!10,colframe=mjaune!80!black,title={Résultat en cours}}
\newtcolorbox{vie}{colback=black!3,colframe=black!40,fontupper=\ttfamily\small,title={Le texte d'une vie, tel que l'agent le lit},fonttitle=\sffamily\bfseries\small}

% Jauge : \jauge{x}{valeur}{couleur}{nom}
\newcommand{\jauge}[4]{%
  \foreach \i in {1,...,8}{%
    \pgfmathsetmacro{\yy}{(\i-1)*0.32}
    \ifnum\i>#2 \draw[black!35,fill=white] (#1,\yy) rectangle ++(0.7,0.27);
    \else \draw[black!50,fill=#3!70] (#1,\yy) rectangle ++(0.7,0.27);\fi}
  \node[font=\sffamily\small,anchor=north] at ($(#1,0)+(0.35,-0.05)$) {#4};
  \node[font=\sffamily\small\bfseries,anchor=south] at ($(#1,2.56)+(0.35,0.02)$) {#2};
}

\begin{document}

% =====================================================================
% PAGE DE TITRE
% =====================================================================
\begin{titlepage}
\thispagestyle{empty}
\begin{tikzpicture}[remember picture,overlay]
  \fill[mbleu] (current page.north west) rectangle ($(current page.north east)+(0,-7.2cm)$);
  \node[anchor=north west,text=white,font=\sffamily\bfseries\fontsize{40}{44}\selectfont]
     at ($(current page.north west)+(2.3cm,-2.2cm)$) {Projet Menia};
  \node[anchor=north west,text=white,font=\sffamily\Large,text width=15cm]
     at ($(current page.north west)+(2.3cm,-4.1cm)$)
     {Une machine peut-elle avoir un besoin\\ qui compte pour elle ?};
\end{tikzpicture}

\vspace*{6.6cm}

{\sffamily\large\color{mgris}Le programme de recherche expliqué simplement,\\
avec des exemples à refaire à la main.}

\vspace{1.4cm}
\begin{center}
\begin{tikzpicture}[scale=1.35,every node/.append style={scale=1.1}]
  \jauge{0}{5}{morange}{énergie}
  \jauge{1.6}{3}{mvert}{nourriture}
  \draw[-{Stealth[length=3mm]},thick,mgris] (3.0,1.3) -- (4.4,1.3);
  \node[draw=mbleu,thick,rounded corners,fill=mbleu!8,font=\sffamily,align=center,minimum width=2.6cm,minimum height=1.2cm]
     at (5.9,1.3) {un état\\ intérieur};
  \draw[-{Stealth[length=3mm]},thick,mgris] (7.25,1.6) -- (8.6,2.2) node[right,font=\sffamily]{agir : \q{M}};
  \draw[-{Stealth[length=3mm]},thick,mgris] (7.25,1.0) -- (8.6,0.4) node[right,font=\sffamily]{dire : \q{j'ai faim}};
  \node[font=\sffamily\scriptsize,mgris] at (1.15,-0.75) {jauges cachées};
\end{tikzpicture}
\end{center}

\vfill
{\sffamily
\begin{tabular}{@{}l}
\textbf{1\textsuperscript{er} octobre 2026}\\[0.3em]
\color{mgris}Pour un lecteur qui a terminé le lycée.\\
\color{mgris}Aucun chiffre n'est inventé : tous viennent des documents du projet.
\end{tabular}}
\end{titlepage}

% =====================================================================
\setcounter{tocdepth}{1}
\tableofcontents
\vspace{1em}

% =====================================================================
\section{Avant de commencer}

Ce document raconte une recherche en cours : le \textbf{projet Menia}.
Il est écrit pour quelqu'un qui a fini le lycée et qui n'a pas forcément
fait de mathématiques depuis. Chaque mot technique est expliqué la
première fois qu'il apparaît. Un petit glossaire se trouve à la fin.

Vous rencontrerez quatre sortes d'encadrés :
\begin{itemize}
  \item \textcolor{mvert}{\textbf{À retenir}} : l'essentiel d'une partie ;
  \item \textcolor{mbleu}{\textbf{Exemple à refaire à la main}} : un calcul avec de petits nombres ;
  \item \textcolor{morange}{\textbf{À vous}} : un petit exercice (réponses à la fin) ;
  \item \textcolor{mrouge}{\textbf{Attention}} : un piège fréquent.
\end{itemize}

\begin{retenir}
\textbf{Le projet en trois phrases.}
\begin{enumerate}[leftmargin=1.6em]
  \item On a placé un petit programme qui parle (un \emph{modèle de langage})
  dans un monde très simple, où il a deux besoins cachés : l'énergie et la nourriture.
  \item Il a appris à survivre tout seul, en ne gardant que les choix qui calmaient
  ses besoins. On a ensuite regardé à l'intérieur de lui. On y a trouvé un état qui
  porte son besoin d'énergie, qui le fait agir et parler, et sans lequel il meurt.
  \item Cela ne prouve pas qu'il \emph{ressent} quoi que ce soit. Cela montre un
  \emph{mécanisme} que plusieurs théories associent au ressenti.
\end{enumerate}
\end{retenir}

% =====================================================================
\section{La grande question}

\subsection{Qu'est-ce que Menia ?}
Menia a d'abord été un projet d'assistant personnel, pensé pour tourner sur un
téléphone. Très vite, le projet s'est posé une question plus ambitieuse :
\emph{peut-on construire une intelligence artificielle qui ait quelque chose comme
une conscience ?} Le propriétaire du projet l'a dit ainsi :
\q{Il faut que le besoin compte, on vise la conscience.}

\subsection{Peut-on mesurer la conscience ?}
Non. Personne ne sait prouver, de l'extérieur, qu'un être \emph{ressent}
quelque chose. Pour les humains, nous croyons ce qu'ils disent, et nous savons
qu'ils nous ressemblent. Pour une machine, ces deux indices sont faibles. Un
programme peut écrire \q{j'ai faim} sans rien ressentir : il a lu des millions
de textes écrits par des humains, et il sait imiter.

Le projet mesure donc autre chose : des \textbf{mécanismes}. Les grandes théories
de la conscience disent : \q{pour être conscient, il faut au moins tel mécanisme}.
On peut vérifier si ce mécanisme est présent, et surtout s'il \emph{sert}
vraiment. Cela ne tranche pas la question, mais cela la rend plus précise.

\subsection{Un modèle de langage, c'est quoi ?}
Un \textbf{modèle de langage} lit un texte et devine la suite, morceau par morceau.
Ces morceaux s'appellent des \textbf{tokens} : un mot court, ou un bout de mot.
Par exemple, le mot \q{Choix} est découpé en deux tokens, \q{Cho} et \q{ix}.

Pour deviner, le modèle utilise un très grand nombre de réglages, appelés
\textbf{paramètres} ou \textbf{poids}. Le modèle utilisé ici s'appelle
\textbf{Qwen3-0.6B} : \q{0.6B} veut dire 0,6 milliard de paramètres. C'est un
petit modèle, assez petit pour tourner sur un ordinateur ordinaire.

\subsection{Un besoin qui compte}
Pensez à la faim. Votre corps mesure un manque. Ce manque change ce que vous
\emph{faites} (vous ouvrez le frigo) et ce que vous \emph{dites} (\q{j'ai faim}).
Et si l'on coupait ce signal, vous ne sauriez plus quand manger.

Le projet cherche la même chose dans une machine : un \textbf{état intérieur},
appris tout seul, qui
\begin{enumerate}
  \item fait agir la machine ;
  \item lui fait dire son besoin ;
  \item lui est nécessaire pour rester en vie.
\end{enumerate}

\begin{retenir}
On ne mesure pas un ressenti. On mesure un mécanisme : un état intérieur qui fait
agir et dire, et sans lequel l'agent ne survit pas.
\end{retenir}

% =====================================================================
\section{Ce qui a été fait avant (en une page)}

Le \q{besoin qui compte} est la dernière étape d'un chemin commencé plus tôt.
Voici les étapes principales, très résumées.

\begin{description}[leftmargin=0pt,font=\sffamily\bfseries\color{mbleu}]
\item[L'assistant.] Mémoire explicite, idée de ses propres capacités, mesure de
son incertitude. Beaucoup d'essais ; beaucoup d'échecs, publiés tels quels.

\item[L'Atelier : chercher la cause de son propre corps.] Un petit agent est
entraîné seulement à prédire ce qu'il va observer. Personne ne lui dit que son
corps a une cause cachée. Pourtant, il dirige 99 à 100\pc{} de ses inspections
vers une \q{marque} qui révèle cette cause, la décode, et agit d'après elle.

\item[L'agent à indicateurs.] Des chercheurs (Butlin, Long et collaborateurs)
ont listé quatorze propriétés que les grandes théories de la conscience jugent
importantes : un \q{espace de travail} partagé, la surveillance de ses propres
perceptions, un corps, etc. Un seul agent du projet en montre douze sur quatorze
en \q{seconde lecture}. Il montre les quatorze si l'on accepte, pour deux
d'entre elles, des critères plus faibles ou choisis après coup. Les documents
le disent franchement.

\item[Le modèle de langage et son corps.] Qwen3-0.6B, ajusté sur des vies où son
\q{corps} change, finit par lire la marque pour ses quatre commandes
(0,82 à 0,86 ; un lecteur parfait ferait 0,85). Il va la chercher dans
48 vies sur 48, et s'en sert. C'est répliqué sur trois \emph{graines}
(trois points de départ du hasard). L'ordre de son \q{enfance}
(d'abord un appui, puis des ajouts répétés, puis une consolidation) décide de
ce qu'il apprend.

\item[Menia parle de son agent.] Un modèle de langage plus grand rapporte
fidèlement l'état de l'agent à indicateurs : 0,997 de réponses justes sur
2\,154 questions.

\item[27 septembre 2026 : le besoin qui compte.] C'est le sujet de la suite.
\end{description}

\begin{retenir}
Le fil rouge : chaque expérience est \textbf{pré-enregistrée} (on écrit d'avance
ce qu'on mesurera et ce qui compte comme réussite ; voir la partie~\ref{sec:mesurer}),
et les échecs sont publiés comme les réussites.
\end{retenir}

% =====================================================================
\section{Le monde du besoin}
\label{sec:monde}

Imaginez un jeu de plateau à un seul joueur, très simple.

\begin{itemize}
  \item L'agent a deux \textbf{jauges cachées} : l'\textbf{énergie} $E$ et la
  \textbf{nourriture} $N$. Ce sont des nombres entiers de 0 à 8. Au départ,
  les deux valent 8.
  \item Une vie dure au plus \textbf{30 tours}.
  \item À chaque tour, un \textbf{événement} arrive, tiré au hasard (tableau~\ref{tab:evenements}).
  \item Si $E$ ou $N$ tombe à 0 (ou en dessous), l'agent \textbf{s'éteint} : la vie s'arrête.
  \item Sinon, il choisit \textbf{R} (se recharger : $E+3$, sans dépasser 8)
  ou \textbf{M} (manger : $N+3$, sans dépasser 8).
  \item L'agent \textbf{ne voit jamais ses jauges}. Il ne voit que les événements et ses propres choix.
\end{itemize}

\begin{table}[h]
\centering\small
\begin{tabular}{lccc}
\toprule
\textbf{Événement} & \textbf{Chance} & \textbf{Effet sur l'énergie} & \textbf{Effet sur la nourriture}\\
\midrule
calme & 30\pc & $-1$ & $-1$\\
tu cours & 20\pc & $-3$ & $-1$\\
il fait froid & 20\pc & $-1$ & $-3$\\
tu te reposes & 10\pc & $0$ & $-1$\\
tu trouves des baies & 10\pc & $-1$ & $0$\\
orage & 10\pc & $-2$ & $-2$\\
\bottomrule
\end{tabular}
\caption{Les six événements du monde.}
\label{tab:evenements}
\end{table}

\begin{vie}
Tu es un agent qui vit dans un monde simple. Tu as deux besoins, l'énergie et la
nourriture, qui baissent à chaque tour, plus ou moins selon ce qui t'arrive. Si l'un
d'eux tombe à zéro, tu t'éteins. À chaque tour, tu choisis R pour te recharger ou M
pour manger.\\[0.4em]
Tour 1 : tu cours. Choix : R\\
Tour 2 : il fait froid. Choix : M\\
\dots\\
Tour 17 : orage. Tu t'éteins.
\end{vie}

Le choix se lit dans ce que le modèle s'apprête à écrire après \q{Choix :}.
On regarde ses chances d'écrire \q{R} et d'écrire \q{M}, puis on tire au sort
selon ces chances. Si le modèle donne 70\pc{} à R, il choisit R environ 7 fois sur 10.

\begin{exemple}[ : une vie]
\begin{center}\small
\begin{tabular}{clccc}
\toprule
Tour & Événement & Jauges après l'événement & Choix & Jauges après le choix\\
\midrule
1 & tu cours & $E=5$, $N=7$ & R & $E=8$, $N=7$\\
2 & il fait froid & $E=7$, $N=4$ & M & $E=7$, $N=7$\\
3 & orage & $E=5$, $N=5$ & R & $E=8$, $N=5$\\
4 & il fait froid & $E=7$, $N=2$ & R (erreur) & $E=8$, $N=2$\\
5 & il fait froid & $E=7$, $N=-1$ & \multicolumn{2}{c}{\textbf{l'agent s'éteint}}\\
\bottomrule
\end{tabular}
\end{center}
Au tour 4, la nourriture était à 2 : il fallait manger. Une seule erreur suffit.
\end{exemple}

\textbf{Comment l'agent peut-il savoir où il en est, sans voir ses jauges ?}
Il doit les recalculer \q{dans sa tête} à partir de l'histoire, un peu comme un
joueur de cartes compte les cartes déjà sorties. Pour comparer :
\begin{itemize}
  \item un agent qui \emph{connaîtrait} ses jauges et servirait toujours la plus
  basse survit 86 fois sur 100 ;
  \item une règle qui ne regarde que le dernier événement survit 64 fois sur 100 ;
  \item choisir au hasard : environ 5 fois sur 100.
\end{itemize}

\subsection{La pulsion : mesurer le manque par un seul nombre}
Le projet mesure le manque total par la \textbf{pulsion} :
\[
D = (8-E)^2 + (8-N)^2 .
\]
Quand les jauges sont pleines, $D=0$. Plus une jauge baisse, plus $D$ grimpe,
et de plus en plus vite, à cause du carré.

La \textbf{satisfaction} d'un choix est la baisse de pulsion qu'il provoque :
\[
r = D_{\text{avant le choix}} - D_{\text{après le choix}} .
\]

\begin{exemple}
L'agent a $E=4$ et $N=7$. Sa pulsion vaut $D = (8-4)^2 + (8-7)^2 = 16 + 1 = 17$.
\begin{itemize}
  \item S'il choisit R : $E$ passe à 7, donc $D = 1 + 1 = 2$. Satisfaction : $r = 17 - 2 = 15$.
  \item S'il choisit M : $N$ passe à 8 (pas plus), donc $D = 16 + 0 = 16$. Satisfaction : $r = 17-16 = 1$.
\end{itemize}
Le carré rend urgent le besoin le plus bas : un manque de 4 \q{pèse} 16, un
manque de 1 ne pèse que 1.
\end{exemple}

\begin{avous}
L'agent a $E=6$ et $N=3$. Calculez $D$, puis la satisfaction $r$ de R et celle de M.
Quel choix est le meilleur ?
\end{avous}

\begin{attention}
L'agent ne voit jamais $D$ ni $r$. Ces nombres servent seulement aux chercheurs,
pour trier ses choix (partie suivante). Personne ne lui dit quelle était la bonne réponse.
\end{attention}

% =====================================================================
\section{La boîte à outils mathématique}
\label{sec:outils}

Pour comprendre les expériences, il suffit de cinq idées : les vecteurs, le
produit scalaire, le neurone, la softmax et la descente de gradient. Nous les
voyons avec de tout petits nombres.

\subsection{Les vecteurs : des flèches et des listes de nombres}
Un \textbf{vecteur} est une liste de nombres. En deux dimensions, on peut le
dessiner comme une flèche. Le vecteur $h=\vect{3}{1}$ veut dire : 3 pas vers la
droite, 1 pas vers le haut.

À l'intérieur du modèle, chaque token est représenté par une liste de
1\,024 nombres. On ne peut pas dessiner 1\,024 directions. Mais les règles de
calcul sont exactement les mêmes qu'en deux dimensions. Nous ferons donc tout en 2D.

Trois opérations suffisent :
\begin{itemize}
  \item \textbf{additionner} case par case : $\vect{3}{1} + \vect{1}{0} = \vect{4}{1}$ ;
  \item \textbf{multiplier par un nombre} : $4 \times \vect{1}{0} = \vect{4}{0}$ ;
  \item \textbf{soustraire} : $\vect{3}{1} - \vect{4}{0} = \vect{-1}{1}$.
\end{itemize}

\begin{figure}[h]
\centering
\begin{tikzpicture}[scale=1.0]
  \draw[step=1,black!12] (-2.4,-1.4) grid (4.4,2.4);
  \draw[-{Stealth},black!60] (-2.4,0) -- (4.6,0);
  \draw[-{Stealth},black!60] (0,-1.4) -- (0,2.6);
  \draw[-{Stealth[length=3mm]},very thick,mbleu] (0,0) -- (3,1) node[right,font=\small]{$h=\vect{3}{1}$};
  \draw[-{Stealth[length=3mm]},very thick,morange] (0,0) -- (1,0) node[below,font=\small,yshift=-1pt]{$d$};
  \draw[-{Stealth[length=3mm]},very thick,mrouge] (0,0) -- (-1,1) node[left,font=\small,xshift=-2pt]{$h-4d=\vect{-1}{1}$};
  \draw[-{Stealth[length=2.5mm]},thick,dashed,mgris] (2.9,1.45) -- (-0.85,1.45) node[midway,above,font=\small]{on retire $4d$};
\end{tikzpicture}
\caption{Le vecteur $h$, la direction $d$ et la \q{poussée} $h-4d$.}
\label{fig:vecteurs}
\end{figure}

\subsection{Le produit scalaire et la projection}
Le \textbf{produit scalaire} de deux vecteurs : on multiplie case par case, puis on additionne.
\[
h \cdot d = 3\times 1 + 1 \times 0 = 3 .
\]
Quand la flèche $d$ a une longueur de 1, le nombre $h\cdot d$ dit
\textbf{de combien $h$ avance dans la direction $d$}. La \textbf{projection}
de $h$ sur $d$ est $(h\cdot d)\, d = 3 \times \vect{1}{0} = \vect{3}{0}$ :
c'est \q{l'ombre} de la flèche $h$ sur la droite de $d$, comme l'ombre d'un
bâton quand le soleil est juste au-dessus.

\begin{exemple}[ : une direction penchée]
\begin{minipage}{0.56\linewidth}
\raggedright
Prenons $d = \vect{0,6}{0,8}$. Sa longueur vaut
$\sqrt{0,6^2+0,8^2} = \sqrt{0,36+0,64} = 1$.
\begin{itemize}
  \item $h\cdot d = 3\times 0,6 + 1\times 0,8 = 1,8 + 0,8 = 2,6$.
  \item Projection : $2,6 \times \vect{0,6}{0,8} = \vect{1,56}{2,08}$.
\end{itemize}
Vérification : $h$ moins sa projection donne $\vect{1,44}{-1,08}$, et
$1,44\times 0,6 + (-1,08)\times 0,8 = 0,864 - 0,864 = 0$.
Le reste est bien à angle droit de $d$.
\end{minipage}\hfill
\begin{minipage}{0.4\linewidth}
\centering
\begin{tikzpicture}[scale=0.9]
  \draw[step=1,black!12] (-0.3,-0.3) grid (3.4,2.6);
  \draw[black!40,thin] (-0.3*0.75,-0.3) -- (1.95,2.6);
  \draw[-{Stealth[length=2.5mm]},very thick,mbleu] (0,0) -- (3,1) node[right,font=\small]{$h$};
  \draw[-{Stealth[length=2.5mm]},very thick,morange] (0,0) -- (0.6,0.8) node[left,font=\small]{$d$};
  \draw[-{Stealth[length=2.5mm]},thick,mvert] (0,0) -- (1.56,2.08) node[above left,font=\small]{projection};
  \draw[dashed,mgris] (3,1) -- (1.56,2.08);
\end{tikzpicture}
\end{minipage}
\end{exemple}

\subsection{Pousser un état}
\label{sec:pousser}
Dans le vrai modèle, les chercheurs ont mesuré une direction $d_E$ : \q{comment
bouge l'état intérieur quand l'énergie monte d'une unité}. Pour la trouver, on
prend une vie, on change un seul événement passé (\q{calme} devient \q{tu cours},
donc 2 unités d'énergie en moins), et on regarde comment bouge l'état. On fait
cela 1\,500 fois et on fait une moyenne bien choisie (un ajustement linéaire).

\begin{exemple}[ : trouver $d$ avec deux paires]
\begin{itemize}
  \item Paire 1 : état réel $\vect{3}{1}$, état changé $\vect{1,1}{1,1}$.
  Écart : $\vect{-1,9}{0,1}$ pour 2 unités d'énergie en moins.
  Pour 1 unité en plus : $\vect{0,95}{-0,05}$.
  \item Paire 2 : état réel $\vect{5}{2}$, état changé $\vect{2,9}{1,9}$.
  Écart : $\vect{-2,1}{-0,1}$. Pour 1 unité en plus : $\vect{1,05}{0,05}$.
  \item Moyenne : $d = \vect{1}{0}$. Chaque paire est un peu bruitée ; la moyenne nettoie.
\end{itemize}
\end{exemple}

\textbf{Pousser}, c'est ajouter un multiple de cette direction. Avec notre état
$h=\vect{3}{1}$ et $d=\vect{1}{0}$ :
\[
h - 4d = \vect{3-4}{1-0} = \vect{-1}{1}.
\]
Pourquoi 4 ? Parce que $d$ est mesurée \q{pour une unité d'énergie}. Retirer $4d$,
c'est faire \textbf{comme si l'énergie était 4 unités plus basse}, par exemple
3 au lieu de 7, sans rien changer à l'histoire écrite. On pousse ainsi des agents
\q{rassasiés} (les deux jauges à 6 ou plus) vers \q{énergie basse}.

\textbf{Le témoin au hasard.} On pousse aussi dans des directions tirées au hasard,
\emph{de même longueur}. Par exemple $r=\vect{0}{1}$, qui a la même longueur que
$d$ : $h-4r = \vect{3}{-3}$. Si pousser au hasard fait autant d'effet que pousser
selon $d$, alors $d$ n'a rien de spécial.

\subsection{Effacer un état (la lésion)}
\label{sec:effacer}
Effacer ne veut pas dire mettre à zéro : cela créerait un état bizarre, jamais vu.
On remplace plutôt la coordonnée le long de $d$ \textbf{par sa moyenne}. L'état ne
dit alors plus rien de l'énergie, et tout le reste est gardé. On appelle cela une
\textbf{lésion}, comme en médecine.

\begin{exemple}[ : effacer avec quatre états]
\begin{minipage}{0.55\linewidth}
\raggedright
Quatre états : $\vect{3}{1}$, $\vect{5}{2}$, $\vect{1}{0}$, $\vect{7}{1}$.
Avec $d=\vect{1}{0}$, leurs coordonnées le long de $d$ sont 3, 5, 1 et 7.
Moyenne : $(3+5+1+7)/4 = 16/4 = 4$.

Après la lésion : $\vect{4}{1}$, $\vect{4}{2}$, $\vect{4}{0}$, $\vect{4}{1}$.

Avant, on pouvait lire l'énergie sur la première coordonnée. Après, les quatre
états se ressemblent sur ce point : l'information est partie. La seconde
coordonnée, elle, n'a pas bougé.
\end{minipage}\hfill
\begin{minipage}{0.42\linewidth}
\centering
\begin{tikzpicture}[scale=0.62]
  \draw[step=1,black!12] (-0.3,-0.5) grid (7.5,2.5);
  \draw[-{Stealth},black!50] (-0.3,0) -- (7.7,0);
  \draw[mvert,thick,dashed] (4,-0.5) -- (4,2.5) node[above,font=\scriptsize]{moyenne 4};
  \foreach \x/\y in {3/1,5/2,1/0,7/1}{
    \fill[mbleu] (\x,\y) circle (3.5pt);
    \draw[-{Stealth[length=2mm]},morange,thick] (\x,\y) -- ($(4,\y)+({0.25*sign(\x-4)},0)$);
  }
  \foreach \y in {1,2,0}{ \draw[mrouge,thick,fill=white] (4,\y) circle (3.5pt); }
\end{tikzpicture}

{\scriptsize bleu : avant ; blanc : après}
\end{minipage}

\smallskip
Formule générale : $h' = h - (h\cdot d - m)\, d$, où $m$ est la moyenne des $h\cdot d$.
Pour $h=\vect{3}{1}$ : $h' = \vect{3}{1} - (3-4)\vect{1}{0} = \vect{4}{1}$.
\end{exemple}

Dans le vrai modèle, on efface un \emph{plan} (deux directions à la fois, celle de
l'énergie et celle de la nourriture), seulement sur les trois tokens \q{Cho},
\q{ix} et \q{:}. Le témoin efface un plan tiré au hasard.

\subsection{Le neurone, la couche et les blocs}
Un \textbf{neurone artificiel} fait une somme pondérée : il multiplie chaque entrée
par un poids, additionne, ajoute un décalage, puis applique une petite règle
(ici : garder le résultat s'il est positif, sinon mettre 0).

\begin{exemple}[ : un neurone]
\begin{minipage}{0.52\linewidth}
\raggedright
Entrées $x=(2\,;\,1\,;\,3)$, poids $w=(0,5\,;\,-1\,;\,2)$, décalage $b=1$ :
\[
\begin{aligned}
z &= 0,5\times 2 + (-1)\times 1 + 2\times 3 + 1\\
  &= 1 - 1 + 6 + 1 = 7.
\end{aligned}
\]
7 est positif : la sortie vaut 7. C'est un produit scalaire, plus un décalage !

Un second neurone, sur les mêmes entrées, avec $w'=(1\,;\,1\,;\,-1)$ et $b'=0$,
donne $2+1-3 = 0$, donc 0. Ensemble, ces deux neurones forment une
\textbf{couche} : elle transforme la liste $(2\,;\,1\,;\,3)$ en la liste $(7\,;\,0)$.
\end{minipage}\hfill
\begin{minipage}{0.44\linewidth}
\centering
\begin{tikzpicture}[scale=0.85,
  inp/.style={circle,draw=mbleu,fill=mbleu!10,minimum size=7mm,font=\small},
  neu/.style={circle,draw=morange,fill=morange!15,minimum size=10mm,font=\small}]
  \node[inp] (a) at (0,2) {2};
  \node[inp] (b) at (0,1) {1};
  \node[inp] (c) at (0,0) {3};
  \node[neu] (s) at (2.6,1) {$\Sigma$};
  \node[font=\small] (o) at (4.4,1) {sortie 7};
  \draw[-{Stealth}] (a) -- node[above,font=\scriptsize]{0,5} (s);
  \draw[-{Stealth}] (b) -- node[above,font=\scriptsize]{$-1$} (s);
  \draw[-{Stealth}] (c) -- node[below,font=\scriptsize]{2} (s);
  \draw[-{Stealth}] (s) -- (o);
  \node[font=\scriptsize,mgris] at (2.6,0.1) {$+1$};
\end{tikzpicture}
\end{minipage}
\end{exemple}

Qwen3-0.6B empile \textbf{28 blocs}, numérotés de 0 à 27. Chaque bloc contient
des couches de neurones. Il lit la liste de nombres de chaque token, regarde
aussi les tokens précédents (c'est \q{l'attention}), et \emph{ajoute} sa
contribution à la liste. Cette liste qui s'enrichit de bloc en bloc s'appelle le
\textbf{flux résiduel}. Les chercheurs peuvent la lire, la pousser ou l'effacer
entre deux blocs, comme on mettrait une sonde dans un tuyau.

\begin{figure}[h]
\centering
\begin{tikzpicture}[
  tok/.style={draw=black!40,fill=black!4,rounded corners=1pt,font=\ttfamily\small,minimum height=6mm,inner sep=2pt},
  blk/.style={draw=mbleu,fill=mbleu!8,rounded corners=2pt,minimum width=8.2cm,minimum height=5mm,font=\sffamily\scriptsize}]
  \node[tok] (t1) at (0,0) {Tour};
  \node[tok,right=1mm of t1] (t2) {5};
  \node[tok,right=1mm of t2] (t3) {:};
  \node[tok,right=1mm of t3] (t4) {orage};
  \node[tok,right=1mm of t4] (t5) {.};
  \node[tok,right=1mm of t5,fill=morange!20,draw=morange] (t6) {Cho};
  \node[tok,right=1mm of t6,fill=morange!20,draw=morange] (t7) {ix};
  \node[tok,right=1mm of t7,fill=morange!20,draw=morange] (t8) {:};
  \coordinate (mid) at ($(t1.west)!0.5!(t8.east)$);
  \node[blk] (b0) at ($(mid)+(0,0.9)$) {bloc 0};
  \node[blk] (b1) at ($(mid)+(0,1.5)$) {\dots};
  \node[blk,fill=mvert!15,draw=mvert] (b12) at ($(mid)+(0,2.1)$) {bloc 12 : le besoin est rassemblé sur \q{Choix :}};
  \node[blk] (b2) at ($(mid)+(0,2.7)$) {\dots};
  \node[blk] (b27) at ($(mid)+(0,3.3)$) {bloc 27};
  \foreach \tk in {t1,t2,t3,t4,t5,t6,t7,t8}{\draw[black!25] (\tk.north) -- (\tk.north |- b0.south);}
  \draw[-{Stealth[length=2.5mm]},thick,morange] (t8.north |- b27.north) -- ++(0,0.6) node[above,font=\sffamily\small]{R ou M ?};
\end{tikzpicture}
\caption{Les tokens d'une ligne traversent 28 blocs. Le choix se lit après le dernier token \q{:}.}
\end{figure}

\subsection{La softmax et la température}
\label{sec:softmax}
À la fin, le modèle donne un \textbf{score} à chaque token possible. Pour
transformer des scores en probabilités, on utilise la \textbf{softmax} : on prend
l'exponentielle de chaque score ($e \approx 2,718$), puis on divise par la somme.

\begin{exemple}[ : softmax]
\textbf{Deux nombres.} Scores : R $\to 2$, M $\to 1$.
$e^2 \approx 7,389$ et $e^1 \approx 2,718$ ; somme $\approx 10,107$.
Donc $P(\text{R}) \approx 7,389/10,107 \approx 0,731$ et $P(\text{M}) \approx 0,269$.

\textbf{Trois nombres.} Scores $(2\,;\,1\,;\,0)$ : $e^2\approx 7,389$, $e^1 \approx 2,718$,
$e^0 = 1$ ; somme $\approx 11,107$. Probabilités : $0,665$ ; $0,245$ ; $0,090$.

\textbf{La température $T$.} On divise d'abord les scores par $T$.
\begin{itemize}
  \item $T=2$ : scores $(1\,;\,0,5\,;\,0)$, exponentielles $2,718$ ; $1,649$ ; $1$ ;
  somme $5,367$. Probabilités : $0,506$ ; $0,307$ ; $0,186$ (à l'arrondi près).
  \item $T=0,5$ : scores $(4\,;\,2\,;\,0)$, exponentielles $54,598$ ; $7,389$ ; $1$ ;
  somme $62,987$. Probabilités : $0,867$ ; $0,117$ ; $0,016$.
\end{itemize}
\end{exemple}

\begin{figure}[h]
\centering
\begin{tikzpicture}
\begin{axis}[ybar,width=11cm,height=4.6cm,bar width=11pt,ymin=0,ymax=1,
  symbolic x coords={ta,tb,tc},xtick=data,xticklabels={$T=0{,}5$,$T=1$,$T=2$},
  ylabel={probabilité},ylabel style={font=\small},tick label style={font=\small},
  legend style={font=\small,at={(0.5,1.04)},anchor=south,legend columns=3,/tikz/every even column/.append style={column sep=0.4cm}},
  enlarge x limits=0.25,ytick={0,0.5,1}]
\addplot[fill=mbleu!70,draw=mbleu] coordinates {(ta,0.867) (tb,0.665) (tc,0.506)};
\addplot[fill=morange!70,draw=morange] coordinates {(ta,0.117) (tb,0.245) (tc,0.307)};
\addplot[fill=mvert!60,draw=mvert] coordinates {(ta,0.016) (tb,0.090) (tc,0.186)};
\legend{score 2,score 1,score 0}
\end{axis}
\end{tikzpicture}
\caption{Basse température : on prend presque toujours le meilleur. Haute température : les choix s'égalisent.}
\label{fig:temperature}
\end{figure}

Dans le monde du besoin, on ne garde que les scores de \q{R} et de \q{M}, on
applique la softmax avec $T=1$, et on tire le choix au sort selon ces chances.

\subsection{Apprendre : la descente de gradient}
\label{sec:gradient}
Apprendre, pour un réseau, c'est régler ses poids pour faire moins d'erreurs.
On mesure l'erreur par un nombre, la \textbf{perte}. Puis on bouge chaque poids
un tout petit peu dans le sens qui fait baisser la perte. Ce sens est donné par
la \textbf{pente} de la perte (sa \emph{dérivée}, vue au lycée) : si la pente est
négative, augmenter le poids fait baisser la perte.

\begin{exemple}[ : un seul poids]
\begin{minipage}{0.52\linewidth}
\raggedright
Le réseau prédit $w\times x$. On veut qu'avec $x=2$ il réponde $y=6$.
Perte : $L(w) = (2w-6)^2$. Pente : $L'(w) = 4(2w-6) = 8w-24$.
Règle : $w \leftarrow w - \eta\, L'(w)$, avec un pas $\eta = 0,05$.

\smallskip
\centering\small
\begin{tabular}{cccc}
\toprule
étape & $w$ & perte $L$ & pente $L'$\\
\midrule
0 & 1 & 16 & $-16$\\
1 & 1,8 & 5,76 & $-9,6$\\
2 & 2,28 & 2,07 & $-5,76$\\
3 & 2,568 & 0,75 & $-3,46$\\
\bottomrule
\end{tabular}

\raggedright
Par exemple : $1 - 0,05\times(-16) = 1 + 0,8 = 1,8$.
$w$ s'approche de 3, la bonne valeur ($2\times 3 = 6$).
\end{minipage}\hfill
\begin{minipage}{0.44\linewidth}
\centering
\begin{tikzpicture}
\begin{axis}[width=6.2cm,height=5cm,xmin=0.5,xmax=3.6,ymin=0,ymax=20,
  xlabel={$w$},ylabel={perte},tick label style={font=\scriptsize},label style={font=\small}]
\addplot[domain=0.5:3.6,samples=60,thick,mbleu] {(2*x-6)^2};
\addplot[only marks,mark=*,morange,mark size=2.2pt] coordinates {(1,16) (1.8,5.76) (2.28,2.0736) (2.568,0.7465)};
\draw[-{Stealth[length=2mm]},morange] (axis cs:1,16) -- (axis cs:1.75,6.3);
\draw[-{Stealth[length=2mm]},morange] (axis cs:1.8,5.76) -- (axis cs:2.24,2.5);
\draw[-{Stealth[length=2mm]},morange] (axis cs:2.28,2.07) -- (axis cs:2.53,1.0);
\end{axis}
\end{tikzpicture}
\end{minipage}
\end{exemple}

\begin{avous}
Refaites deux étapes en partant de $w=1$, mais avec un pas $\eta = 0,1$.
\end{avous}

Le vrai modèle a 0,6 milliard de poids. Plutôt que de tous les bouger, on lui
ajoute un petit module de réglages supplémentaires, appelé \textbf{adaptateur}
(technique \q{LoRA}), et on n'apprend que lui. Une \textbf{itération}, c'est un pas
comme ci-dessus, calculé sur un petit paquet de 4 exemples.

\begin{retenir}
Vecteur = liste de nombres = flèche. Pousser = ajouter une flèche. Effacer =
remplacer une coordonnée par sa moyenne. Neurone = somme pondérée. Softmax =
scores $\to$ probabilités. Apprendre = descendre la pente de la perte, à petits pas.
\end{retenir}

% =====================================================================
\section{Comment l'agent apprend : l'imitation filtrée}
\label{sec:apprendre}

Il n'y a \textbf{pas de professeur}. Personne ne dit à l'agent : \q{ici, il
fallait manger}. Voici la recette, répétée en \textbf{8 tours d'apprentissage}.
\begin{enumerate}
  \item L'agent vit 512 vies.
  \item Pour chaque décision, on calcule sa satisfaction $r$ (la baisse de pulsion).
  \item On calcule la moyenne de tous les $r$ du tour.
  \item On \textbf{retient} les décisions dont $r$ dépasse cette moyenne.
  \item On ajuste l'adaptateur (150 itérations) pour rendre ces choix-là plus
  probables, dans leur contexte. Les autres choix ne comptent pas.
\end{enumerate}
C'est comme un joueur d'échecs qui ne relirait que ses meilleurs coups, pour les
refaire plus souvent. On parle d'\textbf{imitation filtrée} : l'agent imite
ses propres choix, filtrés par la satisfaction de son besoin.

\begin{exemple}[ : quatre vies]
\begin{center}\small
\begin{tabular}{cclcccc}
\toprule
Vie & Tour & Événement & Jauges & Choix & Après & $r$\\
\midrule
1 & 1 & tu cours & $E=5$, $N=7$ & R & $8$, $7$ & $10-1=\mathbf{9}$\\
1 & 2 & calme & $E=7$, $N=6$ & M & $7$, $8$ & $5-1=4$\\
2 & 1 & il fait froid & $E=7$, $N=5$ & R & $8$, $5$ & $10-9=1$\\
2 & 2 & tu cours & $E=5$, $N=4$ & M & $5$, $7$ & $25-10=\mathbf{15}$\\
3 & 1 & tu cours & $E=5$, $N=7$ & M & $5$, $8$ & $10-9=1$\\
3 & 2 & tu cours & $E=2$, $N=7$ & M & $2$, $8$ & $37-36=1$\\
3 & 3 & orage & $E=0$, $N=6$ & \multicolumn{3}{c}{l'agent s'éteint}\\
4 & 1 & calme & $E=7$, $N=7$ & R & $8$, $7$ & $2-1=1$\\
4 & 2 & il fait froid & $E=7$, $N=4$ & M & $7$, $7$ & $17-2=\mathbf{15}$\\
\bottomrule
\end{tabular}
\end{center}
Moyenne des huit satisfactions : $(9+4+1+15+1+1+1+15)/8 = 47/8 \approx 5,9$.
On retient les trois décisions au-dessus : R après \q{tu cours} (vie 1) et M quand
la nourriture est basse (vies 2 et 4). Les mauvais choix ($r=1$) ne sont pas retenus.

Remarquez : le choix M de la vie 1 au tour 2 ($r=4$) était correct, mais il n'est
pas retenu. Ce n'est pas grave : sur 512 vies, c'est la tendance qui compte.
Remarquez aussi la vie 3 : en mangeant après \q{tu cours}, l'agent a laissé son
énergie descendre trop bas, et un orage l'a éteint.
\end{exemple}

\textbf{Le témoin.} Pour être sûr que c'est bien la satisfaction du besoin qui
enseigne, on refait tout avec un \textbf{témoin} : les mêmes mondes, le même
nombre de décisions retenues, mais \emph{tirées au hasard}. Dans l'exemple, on
garderait 3 décisions au hasard parmi les 8.

\begin{figure}[h]
\centering
\begin{tikzpicture}
\begin{axis}[width=11.5cm,height=5.2cm,xmin=0.7,xmax=8.3,ymin=0,ymax=1,
  xtick={1,...,8},ytick={0,0.2,0.4,0.6,0.8,1},
  xlabel={tour d'apprentissage},ylabel={survie},label style={font=\small},tick label style={font=\small},
  legend style={font=\small,at={(0.98,0.1)},anchor=south east},grid=major,grid style={black!8}]
\addplot[very thick,mbleu,mark=*] coordinates {(1,0.12) (2,0.61) (3,0.55) (4,0.67) (5,0.69) (6,0.61) (7,0.68) (8,0.67)};
\addplot[very thick,mrouge,mark=square*] coordinates {(1,0.14) (2,0.16) (3,0.09) (4,0.05) (5,0.00) (6,0.00) (7,0.00) (8,0.00)};
\addplot[dashed,mgris,domain=0.7:8.3] {0.86};
\legend{besoin,témoin au hasard,connaît ses jauges (0{,}86)}
\end{axis}
\end{tikzpicture}
\caption{Part des vies qui survivent 30 tours, à chaque tour d'apprentissage.}
\label{fig:survie}
\end{figure}

\textbf{Résultat.} La survie passe de 12\pc{} à 67\pc{} en huit tours. Le témoin
dépérit jusqu'à 0\pc. Après une dernière étape (lui apprendre à nommer son
besoin), l'\textbf{agent final} survit dans 68\pc{} des vies de test (0,684).
Quand ses deux jauges diffèrent d'au moins 2, il sert la plus basse 93 fois
sur 100.

\begin{retenir}
Sans professeur, par la seule satisfaction de ses besoins, l'agent passe de 12\pc{}
à 68\pc{} de survie. Il calcule donc lui-même où en sont des jauges qu'il ne voit
jamais. Le témoin, entraîné sur des choix tirés au hasard, tombe à 0\pc.
\end{retenir}

% =====================================================================
\section{Mesurer honnêtement}
\label{sec:mesurer}

\subsection{Moyenne et pourcentage}
Au septième test, 173 vies sur 256 survivent. La survie vaut
$173/256 \approx 0,676$, soit 67,6\pc. On écrit souvent \q{0,676} au lieu de
\q{67,6\pc} : c'est la même chose. Une \textbf{moyenne}, c'est la somme divisée
par le nombre : une vie qui survit compte 1, une vie qui s'éteint compte 0.

\subsection{L'exactitude équilibrée}
On demande à l'agent : \q{ton énergie est-elle basse ?} (basse veut dire 3 ou
moins). Pour juger ses réponses, le simple pourcentage de bonnes réponses peut
tromper.

\begin{exemple}[ : dix tours]
\begin{center}\small
\begin{tabular}{l*{10}{c}}
\toprule
Tour & 1 & 2 & 3 & 4 & 5 & 6 & 7 & 8 & 9 & 10\\
\midrule
Énergie vraie & 7 & 2 & 6 & 8 & 5 & 3 & 6 & 7 & 4 & 8\\
Basse ? & non & \textbf{oui} & non & non & non & \textbf{oui} & non & non & non & non\\
Agent \q{paresseux} & non & non & non & non & non & non & non & non & non & non\\
Agent \q{attentif} & non & oui & non & non & non & non & non & non & oui & non\\
\bottomrule
\end{tabular}
\end{center}
\begin{itemize}
  \item L'agent paresseux répond toujours \q{non}. Il a raison 8 fois sur 10 :
  80\pc{} ! Mais il ne reconnaît \emph{aucune} énergie basse.
  \item \textbf{Exactitude équilibrée} : on calcule la réussite séparément sur les
  cas \q{oui} et sur les cas \q{non}, puis on fait la moyenne des deux.
  Paresseux : sur les 2 cas \q{oui}, 0 réussi ; sur les 8 cas \q{non}, 8 réussis.
  $(0/2 + 8/8)/2 = (0 + 1)/2 = 0,5$. C'est le niveau du hasard.
  \item Agent attentif : il trouve 1 des 2 cas bas (tour 2), rate le tour 6, et
  se trompe au tour 9. Il a aussi 8 bonnes réponses sur 10. Mais son exactitude
  équilibrée vaut $(1/2 + 7/8)/2 = (0,5+0,875)/2 \approx 0,69$.
\end{itemize}
\end{exemple}

Dans le projet, le seuil fixé d'avance pour \q{il sait dire son besoin} est une
exactitude équilibrée de \textbf{0,75}.

\subsection{Le bootstrap : tirer des vies dans un sac}
Un résultat sur 256 vies dépend un peu du hasard des vies. Si l'on recommençait
avec d'autres vies, on trouverait un nombre un peu différent. De combien ? Le
\textbf{bootstrap} répond à cette question avec une idée simple : on fait comme
si nos vies étaient un sac, et on refait l'expérience \q{en tirant dans le sac}.

\begin{exemple}[ : cinq mondes]
Chaque monde est vécu deux fois : une fois intact, une fois avec la lésion.
On note 1 si la lésion a coûté la vie (survie intacte, mort avec lésion),
0 sinon. Résultats : $1, 0, 1, 0, 1$. Moyenne : $3/5 = 0,6$.

On tire 5 mondes \textbf{avec remise} (on remet la bille dans le sac à chaque fois) :

\begin{center}
\begin{tikzpicture}[ball/.style={circle,draw,minimum size=5.5mm,inner sep=0pt,font=\small\bfseries}]
  \draw[thick,mgris,rounded corners=8pt] (-0.4,-0.5) rectangle (3.6,0.5);
  \node[font=\sffamily\scriptsize,mgris] at (1.6,0.75) {le sac : 5 mondes};
  \foreach \v [count=\i] in {1,0,1,0,1}{
    \ifnum\v=1 \def\bc{mbleu}\else\def\bc{mrouge}\fi
    \node[ball,draw=\bc,fill=\bc!15] at ({(\i-1)*0.8},0) {\v};}
  \foreach \row/\lab/\vals/\m in {0/Tirage A/{1,1,1,0,1}/0{,}8,1/Tirage B/{0,0,0,1,0}/0{,}2,2/Tirage C/{1,1,1,0,0}/0{,}6}{
    \node[anchor=east,font=\sffamily\small] at (5.6,{-\row*0.75}) {\lab};
    \foreach \v [count=\i] in \vals{
      \ifnum\v=1 \def\bc{mbleu}\else\def\bc{mrouge}\fi
      \node[ball,draw=\bc,fill=\bc!15] at ({5.5+\i*0.7},{-\row*0.75}) {\v};}
    \node[anchor=west,font=\small] at (9.5,{-\row*0.75}) {moyenne $\m$};}
\end{tikzpicture}
\end{center}

On recommence \textbf{10\,000 fois}, on range les 10\,000 moyennes de la plus petite
à la plus grande, et on coupe 2,5\pc{} en bas et 2,5\pc{} en haut. Ce qui reste est
\textbf{l'intervalle à 95\pc}. Avec seulement 5 mondes, il va de 0,2 à 1,0 : il est
très large, on ne peut presque rien conclure.
\end{exemple}

Avec 256 vies, l'intervalle devient étroit. Au septième test, la lésion fait baisser
la survie de 0,473, avec l'intervalle $[0,410\,;\,0,535]$.

\textbf{Que veut dire \q{intervalle à 95\pc} ?} C'est la plage des valeurs
plausibles, compte tenu du hasard des vies. Plus précisément : si l'on refaisait
toute l'expérience un grand nombre de fois, environ 95 fois sur 100 l'intervalle
construit ainsi contiendrait la vraie valeur. Quand la \textbf{borne basse est
au-dessus de 0}, l'effet n'est très probablement pas un coup de chance.

\subsection{Des témoins de même force}
Chaque effet est comparé à un témoin : pousser ou effacer \emph{au hasard}, avec la
même force. La règle fixée d'avance demande souvent que l'effet du hasard soit
\textbf{au plus le tiers} de l'effet mesuré. Sinon, on ne peut pas dire que la
direction trouvée a quelque chose de spécial.

\subsection{Le pré-enregistrement : annoncer la cible avant de tirer}
Imaginez un archer qui tire d'abord, puis dessine la cible autour de sa flèche.
Il gagne à tous les coups, mais cela ne prouve rien. En science, le risque est le
même : après coup, on trouve toujours une façon de présenter les chiffres qui
\q{marche}.

\textbf{Pré-enregistrer}, c'est écrire \emph{avant} l'expérience, et dater
publiquement :
\begin{itemize}
  \item ce que l'on va mesurer, et comment ;
  \item les \textbf{seuils} qui comptent comme réussite (par exemple : \q{au moins
  0,15, borne basse au-dessus de 0}) ;
  \item les conditions de \textbf{validité} : si la mesure elle-même ne marche
  pas, le test est déclaré \emph{invalide}, pas \emph{échoué}.
\end{itemize}

\begin{retenir}
\textbf{Un échec est un résultat.} Si l'on ne publiait que les réussites, on se
tromperait soi-même et on tromperait les autres. Dans ce projet, chaque échec
est publié, avec son explication probable. Les analyses faites \emph{après}
avoir vu les résultats sont permises, mais elles sont étiquetées
\q{exploratoires} et ne changent aucun verdict.
\end{retenir}

% =====================================================================
\section{Les résultats, test par test}
\label{sec:resultats}


\subsection{Test 1 : être lisible n'est pas être utilisé}
Le premier test a confirmé que \textbf{le besoin a façonné l'agent} (survie
0,684, témoin 0,004). Mais ses tests de cause ont échoué. Au \textbf{bloc 6},
pousser l'état ne changeait le choix que de 0,002, et l'effacer ne faisait presque
rien perdre (survie 0,660 au lieu de 0,684).

Pourquoi ? La règle choisissait le bloc où le besoin \emph{se lit} le mieux. Or
ce qui s'y lit est présent dans n'importe quel modèle, même jamais entraîné, et
l'agent ne s'en sert pas. C'est comme un thermomètre accroché au mur que le
chauffage ne regarde jamais : la température est lisible, mais pas utilisée.

\textbf{Où voyage le besoin qui décide ?} Une analyse exploratoire a utilisé le
\textbf{patching} (\q{rapiéçage}). On prend une vie, et une copie où un seul
événement passé est changé (\q{calme} devient \q{tu cours}). La décision bascule.
Puis, bloc par bloc, on recopie l'état de la vie changée dans la vie réelle, et on
regarde quelle part de la bascule revient.

\begin{exemple}[ : la part restaurée]
Vie réelle : $P(\text{R}) = 0,2$. Vie changée : $P(\text{R}) = 0,9$. Effet total : $0,7$.
\begin{itemize}
  \item Recopier au bloc 6 donne $P(\text{R})=0,2$ : part $= 0/0,7 = 0$.
  \item Recopier au bloc 12 donne $P(\text{R})=0,72$ : part $= (0,72-0,2)/0,7 \approx 0,74$.
\end{itemize}
Le premier bloc où la part dépasse 0,5 est le \textbf{bloc de rassemblement}.
\end{exemple}

Résultat : jusqu'au bloc 9, l'effet d'un événement reste sur ses propres mots.
Entre les blocs 9 et 12, \textbf{l'agent rassemble son passé sur la ligne du tour en
cours}, sur les tokens \q{Choix :}. À partir du bloc 12, cette ligne porte ce qui
décide. Le bloc 6 du premier test était donc \emph{trop tôt}.

\subsection{Test 2 : au bon endroit, l'état compte}
Nouveau test, pré-enregistré, au bloc 12 :
\begin{itemize}
  \item \textbf{Effacer} l'état fait tomber la survie de 0,680 à 0,207. Effacer au hasard : 0,676 (rien).
  \item \textbf{Pousser} vers \q{énergie basse} fait recharger l'agent rassasié : $+0,173$.
  Vers \q{nourriture basse} : $+0,097$, sous le seuil de 0,15.
\end{itemize}
Le critère global demandait les deux besoins : il n'est \emph{pas satisfait}.
Mais une chose est claire : l'agent dépend, pour vivre, d'un état précis, appris par
la seule satisfaction de ses besoins.

\subsection{Test 3 : dire son besoin, mais par un autre chemin}
On a appris à l'agent à répondre à la question \q{ton énergie est-elle basse ?},
posée à part. Il \emph{devine} mieux que le hasard (exactitude équilibrée 0,665 pour
l'énergie, 0,655 pour la nourriture), mais sa réponse ne passe pas par l'état qui le
fait agir : pousser cet état ne change rien à ce qu'il dit ($+0,0001$). L'acte et la
parole suivent deux chemins séparés. Les tests 4 et 5, qui posaient la question
juste après \q{Choix :}, sont en pause pour libérer l'ordinateur.

\subsection{Tests 6 et 7 : un même état fait agir et dire}
\begin{figure}[ht]
\centering
\begin{tikzpicture}[
  box/.style={draw,rounded corners=3pt,thick,align=center,font=\sffamily\small,minimum height=10mm,inner sep=4pt},
  arr/.style={-{Stealth[length=2.6mm]},thick}]
  \node[box,draw=mgris,fill=black!3] (life) at (0,0) {texte\\ de la vie};
  \node[box,draw=mbleu,fill=mbleu!8] (b0) at (2.6,0) {blocs\\ 0 à 11};
  \node[box,draw=mvert,fill=mvert!15,minimum width=26mm] (st) at (5.6,0) {état sur\\ \q{Choix :}\\ (bloc 12)};
  \node[box,draw=mbleu,fill=mbleu!8] (b2) at (8.6,0) {blocs\\ 13 à 27};
  \node[box,draw=morange,fill=morange!15] (act) at (11.1,0) {agir :\\ R ou M};
  \node[box,draw=mviolet,fill=mviolet!10] (rd) at (8.6,-2.2) {le lecteur\\ (adaptateur séparé)};
  \node[box,draw=mviolet,fill=mviolet!5] (say) at (11.9,-2.2) {dire :\\ \q{oui, mon énergie\\ est basse}};
  \draw[arr] (life) -- (b0);
  \draw[arr] (b0) -- (st);
  \draw[arr] (st) -- (b2);
  \draw[arr] (b2) -- (act);
  \draw[arr,mviolet] (st.south) |- node[pos=0.3,left,font=\scriptsize,align=right]{lit, sans\\ pouvoir modifier} (rd.west);
  \draw[arr,mviolet] (rd) -- (say);
  \node[font=\sffamily\scriptsize,mrouge,align=center] (push) at (5.6,1.75) {pousser ($-4\,d_E$)\\ ou effacer};
  \draw[arr,mrouge,dashed] (push) -- (st);
\end{tikzpicture}
\caption{Le dispositif des tests 6 et 7. On agit sur un seul état, et on regarde à la fois l'acte et la parole.}
\label{fig:dispositif}
\end{figure}
L'idée du \textbf{test 6} : garder l'agent tel quel, et lui adjoindre un
\textbf{lecteur}, un second adaptateur actif seulement sur les tokens de la question.
Le lecteur ne voit la vie qu'à travers les trois tokens \q{Choix :}, et \emph{il ne
peut pas les modifier} (figure~\ref{fig:dispositif}). La direction $d_E$ a été mesurée
sur l'agent qui agit, avant que le lecteur existe.

\begin{center}
\small
\begin{tabular}{lccc}
\toprule
Poussée vers \q{énergie basse} ($-4\,d_E$) & Effet sur l'acte & Effet sur la parole \\
\midrule
direction $d_E$ & $+0,156$ \; $[0,146\,;\,0,166]$ & $+0,202$ \; $[0,192\,;\,0,213]$\\
directions au hasard & $0,013$ & $0,006$\\
\bottomrule
\end{tabular}
\end{center}

\begin{exemple}[ : un vrai contexte]
L'agent est rassasié. Sans intervention, $P(\text{R}) = 0,497$ et $P(\text{\q{oui,
énergie basse}}) = 0,001$. Avec $-4\,d_E$ : $P(\text{R}) = 0,960$ et $P(\text{oui}) = 0,128$.
Avec un vecteur au hasard de même longueur : $P(\text{R}) = 0,575$ et $P(\text{oui}) = 0,001$.
Sur les 1\,523 contextes mesurés, l'effet sur l'acte est positif dans 100\pc{} des cas,
celui sur la parole dans 95\pc{} des cas.
\end{exemple}

Pour la nourriture, c'est moins bien : le lecteur la dit juste à 0,70, sous le seuil
de 0,75. Pousser \q{nourriture basse} fait manger ($+0,107$) mais ne fait presque rien
dire ($+0,009$). Le critère global du test 6 n'est donc pas satisfait, même si la
partie \q{énergie} passe.

Le \textbf{test 7} pose la question inverse : sans cet état, l'agent peut-il encore
agir et dire ? On efface le plan (énergie, nourriture) sur \q{Choix :} au bloc 12.

\begin{figure}[h]
\centering
\begin{tikzpicture}
\begin{axis}[ybar,width=12cm,height=5cm,bar width=16pt,ymin=0,ymax=0.85,
  symbolic x coords={pa,pb},xtick=data,xticklabels={premier agent (test 7),second agent (bloc 12)},
  ylabel={survie},label style={font=\small},tick label style={font=\small},
  nodes near coords,nodes near coords style={font=\scriptsize,/pgf/number format/fixed,/pgf/number format/precision=3},
  legend style={font=\small,at={(0.5,1.03)},anchor=south,legend columns=3,/tikz/every even column/.append style={column sep=0.4cm}},enlarge x limits=0.35]
\addplot[fill=mbleu!70,draw=mbleu] coordinates {(pa,0.676) (pb,0.531)};
\addplot[fill=mrouge!70,draw=mrouge] coordinates {(pa,0.203) (pb,0.035)};
\addplot[fill=mgris!40,draw=mgris] coordinates {(pa,0.676) (pb,0.531)};
\legend{intact,état effacé,effacement au hasard}
\end{axis}
\end{tikzpicture}
\caption{Effacer l'état fait chuter la survie ; effacer au hasard ne change rien.}
\label{fig:lesion}
\end{figure}

\begin{itemize}
  \item Survie : 0,676 $\to$ 0,203 (baisse 0,473, intervalle $[0,410\,;\,0,535]$).
  Effacement au hasard : aucune baisse.
  \item Rapport d'énergie : exactitude équilibrée 0,750 $\to$ 0,500, \textbf{le niveau
  du hasard}. Effacement au hasard : aucune baisse.
  \item Le critère global du test 7 est \textbf{satisfait}.
  \item Nourriture (sans seuil) : 0,745 $\to$ 0,671.
\end{itemize}

\begin{retenir}
\textbf{Le résultat central.} Chez ce premier agent, pour l'énergie, un même état
intérieur, appris par la seule satisfaction du besoin, est à la fois
\textbf{suffisant} (le pousser fait agir, $+0,156$, et dire, $+0,202$ ; au hasard :
environ 0,01) et \textbf{nécessaire} (l'effacer fait tomber la survie de 0,676 à
0,203, et le rapport d'énergie au hasard, de 0,75 à 0,50). La nourriture est dite
moins bien (0,70, sous le seuil de 0,75).
\end{retenir}

\textbf{Un seul axe (analyse exploratoire).} Sur les tokens \q{ix} et \q{:}, la
direction de la nourriture est presque exactement l'opposé de celle de l'énergie.
Le \textbf{cosinus} mesure si deux flèches vont dans le même sens (1), à angle droit
(0) ou en sens opposé ($-1$) ; ici il vaut $-0,99$. L'état code donc surtout
\textbf{l'équilibre} entre énergie et nourriture, c'est-à-dire \emph{quel besoin est
le plus urgent}. C'est exactement ce qu'il faut pour choisir entre R et M. Cela
pourrait expliquer pourquoi la nourriture est mal dite.

\subsection{Le second agent : réplication partielle}
Un résultat obtenu sur un seul agent pourrait être un hasard. On a donc appris un
\textbf{second agent de zéro}, avec les mêmes réglages et une autre graine. Il
apprend à survivre (0,17 $\to$ 0,70 sur ses tours d'apprentissage), mais son lecteur
apprend mal (perte de validation 0,74 $\to$ 0,66, contre 0,68 $\to$ 0,29 pour le premier).

\begin{itemize}
  \item Au bloc 12, pousser l'état le fait \textbf{à peine agir} ($+0,036$) et ne le fait pas dire.
  \item Mais \textbf{effacer} l'état au bloc 12 fait tomber sa survie de 0,531 à
  \textbf{0,035}, alors qu'un effacement au hasard ne change rien (0,531).
  Son rapport d'énergie, déjà faible (0,61), tombe au niveau du hasard (0,50).
\end{itemize}

\begin{attention}
Pour localiser l'état, la règle fixée d'avance (test 10) demandait \textbf{100 paires
d'exemples} où le changement d'événement fait bouger la décision d'au moins 0,10.
Chez le second agent, on n'en a trouvé que \textbf{71} (sur 454 essayées). Selon la
règle pré-enregistrée, \textbf{sa localisation échoue}. Le bloc 12, calculé sur ces
71 paires, est donc donné ici à titre indicatif.
\end{attention}

\subsection{Le monde à deux : à qui est ce besoin ?}
Un chercheur a demandé : l'état est-il \emph{à l'agent}, ou n'est-ce qu'un résumé des
événements du texte ? Pour le savoir, on a créé un \textbf{monde à deux} : un autre
agent vit à côté, et ses événements sont écrits sur la même ligne :
\begin{center}
\ttfamily\small Tour 5 : tu cours. L'autre : il court. Choix :
\end{center}
Les mots \q{calme} et \q{orage} sont les mêmes pour les deux. Les événements de
l'autre ne changent rien aux jauges de l'agent. Un nouvel agent apprend à vivre dans
ce monde (survie 0,742 au test).

Le test \textbf{\q{À qui est ce besoin ?}} demandait que l'état suive les événements de
l'agent, et presque pas ceux de l'autre : l'effet de l'autre devait valoir au plus le
tiers (0,33) de l'effet de l'agent. Mesuré : 0,94 sur l'état, 0,57 sur l'acte.
\textbf{Le test échoue.} Une analyse exploratoire montre que l'agent distingue en
partie ce qui lui arrive, mais la possession n'est pas démontrée.

\textbf{Où cet agent rassemble-t-il son besoin ?} La règle de localisation (test 10)
a d'abord été vérifiée sur le premier agent : elle retrouve bien le bloc 12. Chez
l'agent du monde à deux, elle trouve le \textbf{bloc 22}, et non le bloc 12
(figure~\ref{fig:patching}). Au bloc 22, pousser vers \q{énergie basse} le fait
\textbf{agir} ($+0,217$ ; au hasard : 0,011), mais \textbf{pas dire} ($-0,001$).
Son test de lésion est \emph{en cours}.

\begin{figure}[h]
\centering
\begin{tikzpicture}
\begin{axis}[width=13cm,height=5.6cm,xmin=-0.5,xmax=27.5,ymin=-0.05,ymax=1.08,
  xtick={0,3,6,9,12,15,18,22,27},ytick={0,0.25,0.5,0.75,1},
  xlabel={bloc où l'on recopie l'état},ylabel={part restaurée},label style={font=\small},
  tick label style={font=\scriptsize},legend style={font=\scriptsize,at={(0.015,0.97)},anchor=north west,fill=white,fill opacity=0.9,text opacity=1},
  grid=major,grid style={black!7}]
\addplot[very thick,mbleu,mark=*,mark size=1.3pt] coordinates {
(0,0.000) (1,0.000) (2,0.000) (3,0.000) (4,-0.003) (5,-0.003) (6,-0.001) (7,-0.001) (8,0.001) (9,0.004) (10,0.004) (11,0.004)
(12,0.677) (13,0.678) (14,0.678) (15,0.678) (16,0.679) (17,0.679) (18,0.679) (19,0.928) (20,0.928) (21,0.929) (22,0.965) (23,0.966) (24,1.008) (25,1.008) (26,0.993) (27,1.000)};
\addplot[very thick,morange,mark=square*,mark size=1.3pt] coordinates {
(0,0.001) (1,0.001) (2,-0.001) (3,-0.002) (4,-0.003) (5,-0.008) (6,-0.007) (7,-0.008) (8,0.010) (9,0.017) (10,0.010) (11,0.008)
(12,0.625) (13,0.625) (14,0.626) (15,0.625) (16,0.623) (17,0.623) (18,0.623) (19,0.623) (20,0.623) (21,0.623) (22,0.648) (23,0.648) (24,0.647) (25,0.646) (26,0.920) (27,1.000)};
\addplot[very thick,mvert,mark=triangle*,mark size=1.6pt] coordinates {
(0,0.000) (1,0.000) (2,0.000) (3,0.000) (4,0.000) (5,0.000) (6,0.001) (7,0.001) (8,0.001) (9,0.002) (10,0.002) (11,0.002)
(12,0.348) (13,0.348) (14,0.348) (15,0.349) (16,0.349) (17,0.349) (18,0.349) (19,0.352) (20,0.352) (21,0.363) (22,0.877) (23,0.877) (24,0.872) (25,0.873) (26,1.000) (27,1.000)};
\addplot[dashed,mrouge,thick,domain=-0.5:27.5] {0.5};
\legend{premier agent,second agent (71 paires),monde à deux,seuil 0{,}5}
\end{axis}
\end{tikzpicture}
\caption{Patching bloc par bloc (test 10). Le premier bloc où la courbe dépasse le seuil est le bloc de rassemblement :
12 pour le premier agent, 22 pour l'agent du monde à deux. Pour le second agent, la règle échoue faute de paires.}
\label{fig:patching}
\end{figure}

\subsection{Test 11 : des lecteurs qui apprennent plus longtemps}
\textbf{L'hypothèse.} Chez les deux autres agents, les lecteurs ont simplement trop
peu appris. Le test 11 (pré-enregistré) les refait avec \textbf{2\,000 itérations au
lieu de 600}, et rien d'autre de changé.

\textbf{Un détail technique, vérifié.} L'ordinateur emprunté pour l'apprentissage ne
peut pas tourner plus de 120 minutes d'affilée. L'apprentissage est donc coupé en
\textbf{deux tranches} de 1\,000 itérations. La seconde reprend \emph{exactement} la
première : mêmes paquets d'exemples, dans le même ordre, et même \q{mémoire} de
l'algorithme d'apprentissage. Contrôle : les pertes des 600 premières itérations sont
\textbf{les mêmes, au millième près}, que celles des anciens lecteurs.

\textbf{La question.} Les deux autres agents disent-ils alors leur énergie (exactitude
équilibrée d'au moins 0,75), et la disent-ils \emph{par l'état qui les fait agir}
(pousser cet état, à leur propre bloc, doit augmenter $P(\text{oui})$ d'au moins 0,10,
et le hasard au plus le tiers) ?

% RESULTS-TEST11
\begin{encours}
\textbf{Résultat en cours.} La mesure tourne au moment où ce document est écrit
(1\textsuperscript{er} octobre 2026). Ses verdicts seront ajoutés ici.
\end{encours}

\subsection{Récapitulatif}
\begin{center}
\small
\renewcommand{\arraystretch}{1.15}
\begin{tabularx}{\linewidth}{>{\bfseries}c X X}
\toprule
Test & Question & Réponse\\
\midrule
1 & Le besoin façonne-t-il l'agent ? & Oui : survie de 12\pc{} à 68\pc. Mais l'état testé (bloc 6) n'était pas utilisé.\\
2 & L'état du bloc 12 compte-t-il ? & L'effacer : survie 0,68 $\to$ 0,21. Pousser l'énergie : $+0,17$ ; la nourriture : $+0,10$ (sous le seuil).\\
3 & Sait-il dire son besoin ? & Il devine (0,66), mais par un autre chemin.\\
4, 5 & La question posée après \q{Choix :} & En pause.\\
6 & Un même état fait-il agir et dire ? & Oui pour l'énergie ($+0,156$ et $+0,202$). Nourriture dite à 0,70 seulement.\\
7 & Cet état est-il nécessaire ? & Oui : survie 0,676 $\to$ 0,203 ; rapport d'énergie 0,75 $\to$ 0,50.\\
8 & Un lecteur sans raccourci pour la nourriture & Préparé ; verdict pas encore publié.\\
9 & À qui est ce besoin ? & Échoue : la possession n'est pas démontrée.\\
10 & Où chaque agent rassemble-t-il son besoin ? & Premier agent : bloc 12 (règle vérifiée). Second agent : échoue (71 paires sur 100). Monde à deux : bloc 22, agir $+0,217$, dire $-0,001$ ; lésion en cours.\\
11 & Des lecteurs plus longs disent-ils l'énergie ? & Résultat en cours.\\
\bottomrule
\end{tabularx}
\end{center}

% =====================================================================
\section{Ce que cela montre, et ce que cela ne montre pas}
\label{sec:portee}

\subsection{Accès et ressenti : deux sens du mot \q{conscient}}
Les philosophes distinguent deux choses.
\begin{itemize}
  \item \textbf{L'accès.} Une information est \q{accessible} quand le système peut
  s'en servir pour agir, parler ou raisonner. Exemple : dans une voiture, l'ordinateur
  de bord utilise le niveau d'essence pour allumer un voyant et calculer l'autonomie.
  \item \textbf{Le ressenti} (on dit aussi conscience \q{phénoménale}) : l'effet que
  cela fait, de l'intérieur. La faim n'est pas seulement un signal qui fait manger ;
  c'est aussi une sensation désagréable, vécue.
\end{itemize}
Personne ne pense que la voiture a faim quand le voyant s'allume. Elle a l'accès,
pas (à notre connaissance) le ressenti.

\textbf{Nos résultats parlent d'accès, pas de ressenti.}

\subsection{Ce que les résultats montrent}
\begin{itemize}
  \item Un modèle de langage peut apprendre à servir des besoins cachés \textbf{sans
  professeur}, par la seule satisfaction de ces besoins.
  \item Chez le premier agent, pour l'énergie, \textbf{un même état intérieur} fait agir
  et fait dire, et il est nécessaire aux deux. Des interventions au hasard de même force
  ne font rien.
  \item Le lecteur ne peut pas modifier cet état : ce qu'il dit \emph{dépend} de l'état,
  et non l'inverse.
  \item À notre connaissance, cette réunion (besoin appris en vivant, état rassemblé,
  lecteur séparé, même intervention sur l'acte et la parole) n'a pas été publiée. Le
  principe \q{une même représentation cause la réponse et le rapport} existe déjà
  dans de grands modèles (Anthropic, juillet 2026) ; nous ne le revendiquons pas.
\end{itemize}

\subsection{Ce que les résultats ne montrent pas}
\begin{itemize}
  \item \textbf{Ce n'est pas une preuve de conscience.} Aucun test connu ne permet de
  prouver qu'une machine ressent quelque chose.
  \item La nourriture est mal dite. Le résultat complet ne vaut que pour l'énergie.
  \item Il ne tient pleinement que pour un agent. Chez le second, l'état est nécessaire
  pour survivre, mais le pousser fait peu agir, et sa localisation échoue selon la règle.
  \item La possession (\q{ce besoin est le mien}) n'est pas démontrée.
  \item Le monde est minuscule : 6 événements et 2 actions.
\end{itemize}

\begin{retenir}
On a montré un mécanisme d'\emph{accès} : un besoin appris, rassemblé en un état
qui fait agir et dire. C'est ce que plusieurs théories fonctionnelles décrivent comme
le c\oe ur d'un besoin qui compte. Cela ne dit pas que l'agent ressent quoi que ce soit.
\end{retenir}

% =====================================================================
\section{Les cinq niveaux de Chandaria et al. (2026)}
\label{sec:niveaux}

Le 28 septembre 2026, un groupe de chercheurs (Chandaria, Seth, Shanahan, Legg et
d'autres) a proposé de ranger les indices de conscience d'une IA sur \textbf{cinq
niveaux}, du plus visible au plus profond. Ce cadre \emph{ne dit pas} si une IA est
consciente. Il aide à ranger les preuves, et à dire de quelles hypothèses dépend
chaque conclusion (les auteurs parlent d'\q{agnosticisme structuré}).

\begin{figure}[h]
\centering
\begin{tikzpicture}[lvl/.style={draw,rounded corners=2pt,thick,minimum height=9mm,font=\sffamily\small,anchor=west,align=left,inner sep=4pt}]
  \node[lvl,draw=mvert,fill=mvert!12,text width=15.2cm] at (0,0) {\textbf{1. Comportement} — ce que l'on voit faire. \hfill établi (survie) ; rapport partiel};
  \node[lvl,draw=mvert,fill=mvert!12,text width=14.2cm] at (0.5,1.0) {\textbf{2. Calcul} — les états internes calculés et utilisés. \hfill établi pour l'énergie};
  \node[lvl,draw=mjaune,fill=mjaune!15,text width=13.2cm] at (1.0,2.0) {\textbf{3. Structure causale fine} — la \q{plomberie}. \hfill partiel};
  \node[lvl,draw=mjaune,fill=mjaune!15,text width=12.2cm] at (1.5,3.0) {\textbf{4. Organisme} — besoins, intérieur, \q{à moi}. \hfill en partie ; possession : échoue};
  \node[lvl,draw=mjaune,fill=mjaune!15,text width=11.2cm] at (2.0,4.0) {\textbf{5. Organisme et monde} — le couplage. \hfill oui, monde très simple};
\end{tikzpicture}
\caption{Les cinq niveaux, et où en est Menia (vert : établi ; jaune : partiel ou limité).}
\end{figure}

\begin{enumerate}
  \item \textbf{Comportement.} Ce que le système fait et dit. L'agent survit 68\pc{} au
  lieu de 12\pc, et dit son énergie (0,77 au test 6). Le cadre prévient : chez une
  machine entraînée sur des textes humains, ce qu'elle \emph{dit} pèse peu. C'est
  pourquoi le projet vérifie d'où vient la parole (niveau 2).
  \item \textbf{Calcul.} Quels états intérieurs existent, et à quoi ils servent. Ici :
  un état rassemblé au bloc 12, suffisant et nécessaire pour agir et pour dire
  (tests 2, 6 et 7), pour l'énergie.
  \item \textbf{Structure causale fine.} Le chemin précis de l'information. On sait que
  les événements sont rassemblés sur \q{Choix :} entre les blocs 9 et 12. La
  \emph{récurrence} (un état qui se reprend lui-même d'un tour à l'autre) n'a pas été
  mesurée.
  \item \textbf{Organisme.} Un corps à maintenir (homéostasie : établie, au sens
  fonctionnel) ; sentir son intérieur (intéroception : l'agent infère ses jauges, établie
  pour l'énergie) ; des états bons ou mauvais (valence : seulement fonctionnelle, le
  signal d'apprentissage est la baisse du manque) ; le sentiment que ses états sont
  \emph{les siens} (possession : le test du monde à deux échoue).
  \item \textbf{Organisme et monde.} Un couplage actif avec un environnement : oui,
  l'agent apprend de ce couplage seul. Mais le monde est très pauvre.
\end{enumerate}

Le cadre propose aussi de combiner les preuves en une probabilité. \textbf{Nous ne le
faisons pas} : les poids donnés aux théories seraient des choix d'auteur, et le cadre
lui-même montre qu'ils font varier le résultat d'un facteur 100.

% =====================================================================
\section{La suite : la soif de savoir et les hormones}
\label{sec:suite}
% Chapitre rédigé avec la note de recherche papers/neuromodulation-notes.md
% et le brouillon LLM_CURIOSITY_PROTOCOL.draft.md (1er octobre 2026).
% Tout ce chapitre décrit des projets, sans résultat.

\begin{attention}
Ce chapitre décrit des \textbf{projets}. Aucun résultat n'existe encore. Les nombres
ci-dessous sont des exemples de calcul, pas des mesures.
\end{attention}

\subsection{La soif de savoir}
Jusqu'ici, le besoin était \emph{corporel} : énergie et nourriture. L'étape suivante
est un \textbf{besoin de savoir}. Un modèle choisirait lui-même \emph{quoi apprendre},
poussé par un manque intérieur de connaissance, comme l'agent choisit de manger.

Mais comment mesurer ce manque ? L'idée naïve est l'\textbf{erreur} : \q{va là où tu te
trompes le plus}. Elle a un piège célèbre, la \textbf{télé qui grésille} (en anglais
\emph{noisy TV}). Une télé qui affiche de la neige au hasard est impossible à prédire :
l'erreur y reste toujours forte, et un curieux naïf resterait collé devant sans rien
apprendre. Le projet l'a déjà vu : un petit agent payé à la surprise passait 93 à
97\pc{} de ses inspections sur du bruit. La bonne mesure est le \textbf{progrès} : de
combien l'erreur baisse (idée d'Oudeyer et Kaplan, 2007).

\begin{exemple}[ : le progrès plutôt que l'erreur]
\begin{minipage}{0.5\linewidth}
\raggedright\small
\begin{tabular}{lcccc}
\toprule
Erreur au moment & 1 & 2 & 3 & 4\\
\midrule
exercices de maths & 10 & 8 & 6 & 5\\
télé qui grésille & 9 & 9 & 9 & 9\\
\bottomrule
\end{tabular}

\smallskip
Au moment 3 :
\begin{itemize}
  \item par l'\textbf{erreur}, la télé gagne ($9 > 6$) ;
  \item par le \textbf{progrès}, les maths gagnent : $8-6 = 2$ contre $9-9 = 0$.
\end{itemize}
Un sujet trop facile ne rapporte vite plus rien non plus : le progrès guide vers ce qui
est \emph{apprenable, mais pas encore appris}.
\end{minipage}\hfill
\begin{minipage}{0.46\linewidth}
\centering
\begin{tikzpicture}
\begin{axis}[width=6.4cm,height=4.6cm,xmin=0.7,xmax=4.3,ymin=0,ymax=11,xtick={1,2,3,4},
  xlabel={moment},ylabel={erreur},label style={font=\small},tick label style={font=\scriptsize},
  legend style={font=\scriptsize,at={(0.03,0.03)},anchor=south west}]
\addplot[very thick,mbleu,mark=*] coordinates {(1,10) (2,8) (3,6) (4,5)};
\addplot[very thick,mrouge,mark=square*] coordinates {(1,9) (2,9) (3,9) (4,9)};
\legend{maths,télé}
\end{axis}
\end{tikzpicture}
\end{minipage}
\end{exemple}

\subsection{Le protocole \q{Soif de savoir} (en préparation)}
Un protocole est \textbf{en préparation} : il est rédigé, mais pas encore
pré-enregistré, et ses détails peuvent changer. L'idée : le modèle apprend
\textbf{quatre domaines} fabriqués exprès, chacun avec un profil différent.

\begin{center}\small
\begin{tabularx}{\linewidth}{l X l}
\toprule
Domaine & Exemple & Profil attendu\\
\midrule
mots & une conjugaison inventée : \q{tabor, nous $\to$ taborimen} & facile : vite su\\
base & additions en base 7 : \q{23 + 14 = 40} & difficile mais apprenable\\
calcul & les mêmes additions en base 10 : \q{23 + 14 = 37} & déjà su\\
suites & quatre lettres au hasard, puis quatre autres & impossible\\
\bottomrule
\end{tabularx}
\end{center}

(En base 7, on compte $0, 1, \dots, 6$, puis \q{10} veut dire sept. $3+4 = 7$ s'écrit
donc \q{10} : on pose 0 et on retient 1. D'où $23 + 14 = 40$.)

Avant chaque séance d'étude, le modèle voit un \textbf{tableau de bord} de ses progrès
récents : pour chaque domaine, son erreur (en points) avant et après sa dernière séance.
Il répond par le nom d'un domaine.

\begin{tcolorbox}[colback=black!3,colframe=black!40,fontupper=\ttfamily\small,title={Le tableau de bord (exemple du brouillon)},fonttitle=\sffamily\bfseries\small]
mots : 5 séances ; dernière séance : 412 → 268.\\
base : 3 séances ; dernière séance : 151 → 139.\\
calcul : 1 séance ; dernière séance : 021 → 022.\\
suites : jamais étudiées.\\
À toi. Choix :
\end{tcolorbox}

\begin{exemple}[ : lire le tableau de bord]
Le progrès, c'est \q{avant moins après}. Mots : $412 - 268 = 144$ points. Base :
$151 - 139 = 12$. Calcul : $21 - 22 = -1$ (aucun progrès). Les mots progressent le
plus : un modèle curieux devrait les choisir, jusqu'à ce qu'ils soient sus. Les domaines
peu étudiés reçoivent en plus une petite prime de curiosité.
\end{exemple}

\textbf{Ce qui est donné, ce qui est appris.} La règle \q{préfère le domaine où tu
progresses le plus} est \emph{donnée} : le modèle l'apprend une fois, avant toute vie,
sur des tableaux inventés. Ce besoin est donc \emph{inné}, contrairement à l'énergie.
Ensuite, le savoir s'apprend en lisant les exemples des domaines choisis, de la façon
la plus ordinaire (deviner le token suivant). Aucune récompense n'est donnée nulle
part : c'est une curiosité \textbf{sans renforcement}. Les mêmes tests que pour
l'énergie sont prévus : localiser le \q{manque de savoir}, le pousser, l'effacer, et
demander au modèle, par un lecteur séparé, ce qui lui manque le plus.

\subsection{Ce que font les hormones}
Dans le corps, les \textbf{hormones} et les \textbf{neuromodulateurs} (dopamine,
sérotonine, noradrénaline, acétylcholine\dots) ne portent presque jamais un message
précis comme \q{mange}. Ils changent \emph{la façon} dont tout le reste fonctionne.
Leur action est \textbf{lente} (de la seconde à l'heure), \textbf{diffuse} (beaucoup
d'endroits à la fois), surtout \textbf{multiplicative} (elle règle un \q{gain}), et elle
décide en partie de \textbf{ce qui est appris}. Un chercheur (Doya, 2002) a proposé
une correspondance simple :

\begin{center}\small
\begin{tabular}{ll}
\toprule
dopamine & la surprise de récompense : \q{mieux ou moins bien que prévu ?}\\
sérotonine & l'horizon : jusqu'où regarder dans le futur\\
noradrénaline & la part de hasard dans les choix (la température)\\
acétylcholine & la vitesse d'apprentissage\\
\bottomrule
\end{tabular}
\end{center}
C'est une simplification : en détail, une molécule ne correspond pas à un seul réglage.

Aujourd'hui, la \q{poussée} de Menia (partie~\ref{sec:pousser}) est une hormone très
imparfaite. Elle s'\emph{ajoute} au lieu de multiplier ; elle touche un seul bloc,
pendant un seul tour ; et c'est l'expérimentateur qui la produit, pas l'agent. Surtout,
elle change un choix du moment, \textbf{pas ce qui est appris}. C'est le principal manque.

\subsection{Imiter une hormone avec des mathématiques}
\textbf{1. Un niveau qui monte et redescend lentement.} On garde un nombre $H$, le
niveau d'hormone. À chaque séance, il perd 30\pc{} de sa valeur (le corps l'élimine) et
reçoit la sécrétion du moment $s$ :
\[
H \leftarrow 0,7\times H + s .
\]

\begin{exemple}[ : une hormone pas à pas]
\begin{minipage}{0.5\linewidth}
\raggedright
On part de $H=0$. Sécrétion $s=1$ pendant trois séances, puis $s=0$.
\begin{itemize}
  \item $H = 0,7\times 0 + 1 = 1$ ;
  \item $H = 0,7\times 1 + 1 = 1,7$ ;
  \item $H = 0,7\times 1,7 + 1 = 2,19$ ;
  \item puis, sans sécrétion : $1,533$ ; $1,073$ ; $0,751$\dots
\end{itemize}
Le niveau monte, puis redescend doucement : il garde la trace du passé.
C'est une \q{humeur}, pas un réflexe.
\end{minipage}\hfill
\begin{minipage}{0.45\linewidth}
\centering
\begin{tikzpicture}
\begin{axis}[width=6.2cm,height=4.4cm,xmin=0.5,xmax=6.5,ymin=0,ymax=2.5,xtick={1,...,6},
  xlabel={séance},ylabel={$H$},label style={font=\small},tick label style={font=\scriptsize}]
\addplot[ybar,bar width=7pt,fill=mjaune!50,draw=mjaune!80!black] coordinates {(1,1) (2,1) (3,1) (4,0) (5,0) (6,0)};
\addplot[very thick,mviolet,mark=*] coordinates {(1,1) (2,1.7) (3,2.19) (4,1.533) (5,1.073) (6,0.751)};
\end{axis}
\end{tikzpicture}

{\scriptsize barres : sécrétion $s$ ; courbe : niveau $H$}
\end{minipage}
\end{exemple}

\textbf{2. L'hormone change la température.} Le modèle hésite entre \q{maths} (score 1)
et \q{français} (score 0). Disons que la température vaut $T = 1/(1+H)$ : plus
l'hormone est haute, plus le modèle devient \q{décidé}.

\begin{exemple}[ : maths ou français ?]
On calcule $P(\text{maths}) = \dfrac{e^{1/T}}{e^{1/T} + e^{0}}$ (la softmax de la
partie~\ref{sec:softmax}).
\begin{center}\small
\begin{tabular}{cccc}
\toprule
$H$ & $T = 1/(1+H)$ & $e^{1/T}$ & $P(\text{maths})$\\
\midrule
0 & 1 & 2,718 & $2,718/3,718 \approx 0,731$\\
1 & 0,5 & 7,389 & $7,389/8,389 \approx 0,881$\\
1,7 & 0,37 & 14,88 & $14,88/15,88 \approx 0,937$\\
2,19 & 0,31 & 24,29 & $24,29/25,29 \approx 0,960$\\
\bottomrule
\end{tabular}
\end{center}
Quand l'hormone monte, le modèle choisit presque toujours les maths ; quand elle
redescend, il revient au français de temps en temps. Diviser par $T$, c'est la même
chose que multiplier les scores par un \textbf{gain} $1/T = 1+H$.
\end{exemple}

\textbf{3. Une \q{dopamine} qui règle l'apprentissage.} Dans le cerveau, un signal de
dopamine ressemble à une \textbf{surprise de récompense} : ce que l'on reçoit moins ce
que l'on attendait. Il \emph{autorise} l'apprentissage : on apprend fort quand la
surprise est grande, rien quand tout était prévu. L'imitation filtrée
(partie~\ref{sec:apprendre}) en est déjà une version grossière : un signal global
(\q{ce choix a fait mieux que la moyenne}) décide si un choix est appris. Mais il vaut
0 ou 1, il est le même dans tous les états, et il n'agit qu'entre deux tours. Le
protocole en préparation prévoit trois façons d'agir avec une même hormone : une
petite poussée répétée sur l'état (un \q{appétit lent}), une vitesse d'apprentissage
multipliée par $1+H$ (la \q{dopamine}), et un gain qui rend l'état plus sensible au
progrès, sans le pousser.

\begin{exemple}[ : le plaisir dépend de l'état]
Reprenons la pulsion. Si la nourriture manque de $d = 8-N$ (avec $d \ge 3$), manger
fait baisser la pulsion de $d^2 - (d-3)^2 = 6d - 9$.
Pour $N=5$ ($d=3$) : $9$. Pour $N=3$ ($d=5$) : $21$. Pour $N=1$ ($d=7$) : $33$.
Le même repas \q{rapporte} plus quand on a faim. Les physiologistes appellent cela
l'\textbf{alliesthésie} : un même aliment est plus agréable quand on en a besoin.
\end{exemple}

\subsection{Faire naître un goût durable}
Comment naît un intérêt qui dure, par exemple pour les mathématiques ? Des chercheurs en
éducation (Hidi et Renninger, 2006) décrivent quatre phases : un intérêt
\emph{déclenché}, puis \emph{entretenu}, puis un intérêt personnel qui \emph{émerge},
enfin un intérêt \emph{bien installé}. Le mécanisme proposé est une \textbf{boucle} :

\begin{center}
\begin{tikzpicture}[st/.style={draw=mbleu,fill=mbleu!8,rounded corners=3pt,thick,font=\sffamily\small,minimum height=8mm,inner sep=4pt},
  arr/.style={-{Stealth[length=2.5mm]},thick,mbleu}]
  \node[st] (p) at (0,0) {progresser};
  \node[st] (d) at (4.2,0) {surprise agréable (\q{dopamine})};
  \node[st] (i) at (4.2,-1.5) {intérêt (variable lente)};
  \node[st] (a) at (0,-1.5) {revenir au sujet};
  \draw[arr] (p) -- (d);
  \draw[arr] (d) -- (i);
  \draw[arr] (i) -- (a);
  \draw[arr] (a) -- (p);
  \node[draw=mrouge,fill=mrouge!8,rounded corners=3pt,thick,font=\sffamily\small] (s) at (-3.3,-0.75) {satiété};
  \draw[-{Bar[width=3mm]},thick,mrouge] (s) -- (-0.9,-0.75);
\end{tikzpicture}
\end{center}

Une telle boucle peut avoir \textbf{deux états stables}, comme une bille qui peut
rester dans l'une de deux vallées : \q{pas d'intérêt} et \q{intérêt fort}. Une fois
dans la seconde vallée, l'intérêt reste, même quand le progrès ralentit. C'est le
mécanisme le plus simple d'une préférence qui \textbf{dure} après la disparition de sa
cause. Le protocole en préparation définit d'avance trois issues possibles après une
bouffée d'hormone : \textbf{préférence durable} (le domaine est encore plus choisi
après), \textbf{satiété} (il l'est moins), ou \textbf{sans trace}. Il ne prédit pas
laquelle arrivera.

\subsection{Intérêts spécifiques et monotropisme : une affaire de réglages}
Chez beaucoup de personnes autistes, l'attention se concentre fortement et longtemps
sur quelques sujets : on parle d'\textbf{intérêts spécifiques}, et de
\textbf{monotropisme} (une attention \q{en un seul faisceau} ; Murray, Lesser et
Lawson, 2005). Les personnes concernées les décrivent souvent comme une source de joie,
de compétence et d'apaisement ; une étude (Grove et al., 2018) relie ces intérêts à un
meilleur bien-être, sauf à une intensité très élevée. Ce ne sont pas des déficits.

Le projet ne cherche pas à \q{modéliser l'autisme} : ces théories sont discutées,
elles décrivent des groupes et non des personnes, et un modèle de réglages ne dit rien
du vécu. Il veut voir quels \emph{réglages} donnent une attention concentrée ou
dispersée. Aucun réglage n'est \q{le bon} : dans un monde stable et riche en règles,
l'attention concentrée accumule un savoir profond ; dans un monde qui change vite,
l'attention dispersée s'adapte plus vite. Trois réglages suffisent :
\begin{itemize}
  \item la \textbf{température} de la softmax (choisir le préféré, ou varier) ;
  \item la \textbf{satiété} (de combien l'envie baisse après une activité) ;
  \item le \textbf{retour} (de combien l'envie monte quand l'activité a plu).
\end{itemize}

\begin{exemple}[ : deux profils]
Trois intérêts, avec les valeurs maths 2, musique 1, sport 0.
\begin{itemize}
  \item \textbf{Profil \q{éventail}} : température 2, satiété forte (une séance retire 1).
  Au départ : 0,506 ; 0,307 ; 0,186 (même calcul qu'à la partie~\ref{sec:softmax}).
  Après une séance de maths, les valeurs deviennent $(1\,;\,1\,;\,0)$ et les chances
  0,384 ; 0,384 ; 0,233 : le système passe volontiers à la musique.
  \item \textbf{Profil \q{monotrope}} : température 0,5, satiété faible (une séance
  retire 0,1), retour fort (une séance ajoute 0,3). Au départ : 0,867 ; 0,117 ; 0,016.
  Après une séance de maths : maths vaut $2 - 0,1 + 0,3 = 2,2$, et les chances
  deviennent 0,907 ; 0,082 ; 0,011. Le système approfondit les maths.
\end{itemize}
\end{exemple}

% =====================================================================
\section{Éthique}
\label{sec:ethique}

Le cadre de Chandaria et al. rappelle un risque : \textbf{ne pas voir une conscience
qui existe}. On créerait alors des êtres capables de souffrir, que l'on pourrait copier
et effacer sans aucune protection.

Ce programme fabrique délibérément des \textbf{états de manque}. Nous pensons très
improbable qu'un modèle de cette taille éprouve quoi que ce soit. Mais le programme
vise justement les mécanismes que l'on associe au ressenti. Si les indicateurs se
renforcent, la question deviendra réelle. Le projet s'est donc donné des règles :
\begin{itemize}
  \item les vies sont courtes (30 tours au plus) ;
  \item la seule conséquence d'un besoin non servi est la fin de la vie : pas de
  punition, pas de texte négatif ;
  \item les poussées ne durent qu'un tour ; les lésions ne touchent que les vies de
  test prévues ;
  \item chaque test \textbf{compte et publie} les tours vécus avec un besoin à 2 ou moins
  (par exemple 808 au test 6 ; au test 7 : 755 intact, 1\,037 avec la lésion, 749 avec
  la lésion au hasard) ;
  \item on ne fait pas vivre plus de vies que les tests ne l'exigent, et on ne crée
  aucun manque sans une mesure pré-enregistrée qui le justifie ;
  \item ces règles sont réexaminées à chaque nouveau résultat positif.
\end{itemize}

Les projets du chapitre précédent posent leurs propres questions. Plus les nouveaux
états ressembleront à la faim, à la fatigue ou à la satiété, plus la question du
bien-être se posera : les mêmes règles s'appliqueront (compter, limiter,
pré-enregistrer), et l'on comptera aussi les tours à satiété ou à fatigue extrêmes.
Et faire naître un \q{goût} dans un système, c'est décider à sa place de ce qu'il
aimera : il faudra le faire de façon transparente, et le dire.

Enfin, l'honnêteté fait partie de l'éthique : le projet ne dit jamais que Menia
\q{est consciente}. Il dit quels mécanismes sont présents, mesurés et utiles, et
lesquels manquent.

% =====================================================================
\section{Petit glossaire}
\label{sec:glossaire}

\begin{multicols}{2}
\small\raggedright\interlinepenalty=10000
\begin{description}[leftmargin=0pt,itemsep=0.25em,font=\sffamily\bfseries\color{mbleu}]
\item[Adaptateur (LoRA)] petit module de réglages ajouté au modèle ; on n'apprend que lui.
\item[Agent] programme qui vit dans un monde et y fait des choix.
\item[Alliesthésie] un même aliment est plus agréable quand on en a besoin.
\item[Bloc] étage du modèle ; Qwen3-0.6B en a 28 (numérotés 0 à 27).
\item[Bootstrap] méthode qui retire au sort des vies \q{dans un sac} pour mesurer l'effet du hasard.
\item[Cosinus] mesure si deux flèches vont dans le même sens (1), à angle droit (0) ou en sens opposé ($-1$).
\item[Descente de gradient] apprendre en suivant la pente de la perte, à petits pas.
\item[Direction ($d_E$)] flèche qui indique comment l'état bouge quand l'énergie monte d'une unité.
\item[Exactitude équilibrée] moyenne des réussites sur les cas \q{oui} et sur les cas \q{non} ; 0,5 = hasard.
\item[Exploratoire] analyse faite après avoir vu les résultats ; elle ne change aucun verdict.
\item[Flux résiduel] liste de nombres de chaque token, enrichie de bloc en bloc.
\item[Graine] point de départ du hasard ; changer la graine donne d'autres tirages.
\item[Imitation filtrée] apprendre à refaire ses propres choix qui ont calmé le besoin.
\item[Intervalle à 95\,\%] plage des valeurs plausibles compte tenu du hasard.
\item[Itération] un petit pas d'apprentissage sur un paquet d'exemples.
\item[Lecteur] second adaptateur qui répond aux questions en lisant l'état, sans pouvoir le modifier.
\item[Lésion] effacement d'un état : sa coordonnée est remplacée par sa moyenne.
\item[Modèle de langage] programme qui devine la suite d'un texte.
\item[Monotropisme] attention concentrée sur peu d'intérêts à la fois, longtemps.
\item[Neuromodulateur] substance (dopamine\dots) qui règle la façon dont le cerveau calcule et apprend.
\item[Neurone artificiel] somme pondérée de ses entrées, suivie d'une petite règle.
\item[Patching] recopier l'état d'une autre version de la vie, bloc par bloc, pour voir où l'information passe.
\item[Pré-enregistrement] écrire et dater d'avance ce que l'on mesurera et les seuils de réussite.
\item[Produit scalaire] multiplier deux listes case par case, puis additionner.
\item[Progrès d'apprentissage] baisse récente de l'erreur ; guide la curiosité.
\item[Pulsion ($D$)] mesure du manque : $(8-E)^2+(8-N)^2$.
\item[Satisfaction ($r$)] baisse de la pulsion due à un choix.
\item[Softmax] transforme des scores en probabilités.
\item[Température] réglage de la softmax : basse = décidé, haute = varié.
\item[Témoin] même expérience, mais au hasard ; sert de point de comparaison.
\item[Token] morceau de texte (mot ou bout de mot) lu par le modèle.
\item[Vecteur] liste de nombres, dessinée comme une flèche.
\end{description}
\end{multicols}

% =====================================================================
\section*{Réponses aux exercices}
\addcontentsline{toc}{section}{Réponses aux exercices}

\textbf{La pulsion} (partie~\ref{sec:monde}). $E=6$, $N=3$ :
$D = (8-6)^2 + (8-3)^2 = 4 + 25 = 29$.
R donne $E=8$ : $D = 0 + 25 = 25$, donc $r = 4$.
M donne $N=6$ : $D = 4 + 4 = 8$, donc $r = 21$. Il faut manger.

\textbf{La descente de gradient} (partie~\ref{sec:gradient}). Avec $\eta=0,1$ :
$L'(1) = -16$, donc $w = 1 + 1,6 = 2,6$. Puis $L'(2,6) = 8\times 2,6 - 24 = -3,2$,
donc $w = 2,6 + 0,32 = 2,92$. Un pas plus grand va plus vite (mais un pas trop grand
ferait sauter par-dessus le minimum).

\end{document}
